Using the same 20-day OHLCV information set in the Chinese stock market, we revisit whether chart-image CNNs outperform numerical representations in return prediction. Under a fair chronological validation protocol, CNNs exhibit substantial predictability but do not consistently dominate strong nonlinear numerical learners. We introduce a temporal numerical benchmark, a 1D-CNN trained on the same 120D OHLCV inputs reshaped as a $20\times 6$ sequence. This model is the strongest standalone learner, with Rank IC 0.0484 and equal-weighted high-minus-low Sharpe 3.21 under a standard arithmetic Sharpe convention, showing that temporal inductive bias explains a large part of the CNN-versus-flattened-model residual. Yet the chart CNN is not fully spanned: its pairwise residual conditional on the 1D-CNN remains positive, and a CNN plus 1D-CNN ensemble reaches Sharpe 3.46. Pairwise temporal residual tests show that chart-CNN content is more sensitive to near-end local path and volume structures and is realized chiefly at short horizons, whereas 1D-CNN content conditional on CNN is more persistent and pressure-sensitive. These findings shift the comparison from seeking a universally superior model to understanding how different nonlinear representations extract complementary trading signals from the same price history.
